% Проверьте наличие в преамбуле:
% \usepackage{tikz}
% \usetikzlibrary{positioning, arrows.meta}
% \usepackage{graphicx}

\begin{frame}{Схема исследования}
	\begin{center}
		% Масштабирование по ширине слайда автоматически растянет обновленную 
		% геометрию блоков, ликвидируя пустые поля сверху и снизу
		\resizebox{\textwidth}{!}{
			\begin{tikzpicture}[
				node distance=0.4cm,
				param/.style={rectangle, rounded corners, draw=Burgundy!50, fill=Burgundy!1!Ivory, thick, text width=2.8cm, align=center, font=\small},
				opt/.style={rectangle, rounded corners, draw=Burgundy!80, fill=Burgundy!5, thick, text width=3.2cm, align=center, font=\small},
				solver/.style={rectangle, draw=DeepBlue!60, fill=DeepBlue!5, thick, rounded corners, text width=5.4cm, align=left, font=\scriptsize},
				fitness/.style={rectangle, draw=DeepBlue!60, fill=DeepBlue!5, thick, rounded corners, text width=3.6cm, align=center, font=\small},
				note/.style={text=DeepBlue!80, align=center, font=\tiny\bfseries},
				arrow/.style={-Stealth, thick, DeepBlue!80}
				]
				
				% --- ЦЕНТРАЛЬНЫЕ БЛОКИ ---
				
				% Входные параметры
				\node (start) [param] {
					\textbf{Входные данные} \\
					\vspace{4pt}
					\scriptsize Вектор свободных параметров схемы \\ $\theta \in [0, 1]^4$
				};
				
				% Оптимизатор
				\node (algo) [opt, right=of start] {
					\textbf{Каскад оптимизации} \\
					\vspace{4pt}
					\scriptsize Итерационный поиск минимума функционала $C(\theta)$:\\ 
					Grid Search $\rightarrow$ \\ PSO $\rightarrow$ Нелдер--Мид
				};
				
				% --- ВЕРТИКАЛЬНЫЙ СТЕК ЗАДАЧ КОШИ (РАЗВЕТВЛЕНИЕ) ---
				
				% Задача 2 (базовая для выравнивания)
				\node (t2) [solver, right=0.9cm of algo, yshift=0.8cm] {
					2. Уравнение Шрёдингера 2D \\ \quad ($h = 0.05$, $T \le 2$, расчет волновой массы)
				};
				
				% Задача 1 (сверху от 2)
				\node (t1) [solver, above=of t2] {
					1. Гармонический осциллятор \\ \quad ($h = 0.1$, $T \le 10$, тест фазового сдвига)
				};
				
				% Задача 3 (снизу от 2)
				\node (t3) [solver, right=0.9cm of algo, yshift=-0.8cm] {
					3. Задача двух тел \\ \quad ($h = 0.1$, $T \le 15$, кеплеровская орбита)
				};
				
				% Задача 4 (снизу от 3)
				\node (t4) [solver, below=of t3] {
					4. Задача $N$ тел (внешние планеты) \\ \quad ($h = 0.2$, $T \le 20$, 24-мерная система)
				};
				
				% --- ПОСЛЕДУЮЩИЕ БЛОКИ ---
				
				% Оценка качества (выравнивается по центру относительно блока задач)
				\node (eval) [fitness, right=0.9cm of t2, yshift=-0.8cm] {
					\textbf{Критерий качества} \\
					\vspace{4pt}
					\scriptsize Вычисление среднего значения евклидовой нормы погрешности траектории $C(\theta)$
				};
				
				% Выходные данные
				\node (stop) [param, right=of eval] {
					\textbf{Выходные данные} \\
					\vspace{4pt}
					\scriptsize Оптимизированные параметры таблицы Бутчера ($\theta^*$)
				};
				
				% --- ТЕКСТОВЫЕ АННОТАЦИИ ОБЪЕМА ВЫЧИСЛЕНИЙ ---
				\node (vol_note) [note, above=0.3cm of t1, text width=6cm] {
					МАТЕМАТИЧЕСКИЙ ОБЪЕМ ВЫЧИСЛЕНИЙ:\\ 
					900 полных интегрирований систем ОДУ\\ 
					на весь цикл поиска PSO для каждого бенчмарка \\
					(30 частиц $\times$ 30 итераций)
				};
				
				% --- СВЯЗИ И СТРЕЛКИ ---
				
				\draw [arrow] (start) -- (algo);
				\draw [arrow] (eval) -- node[above, font=\tiny] {Сошлось} (stop);
				
				% Разветвление стрелок от Оптимизатора к Задачам
				\draw [arrow] (algo.east) -- ++(0.3,0) |- (t1.west);
				\draw [arrow] (algo.east) -- ++(0.3,0) |- (t2.west);
				\draw [arrow] (algo.east) -- ++(0.3,0) |- (t3.west);
				\draw [arrow] (algo.east) -- ++(0.3,0) |- (t4.west);
				
				% Схождение стрелок от Задач к Критерию качества
				\draw [arrow] (t1.east) -- ++(0.3,0) |- (eval.west);
				\draw [arrow] (t2.east) -- ++(0.3,0) |- (eval.west);
				\draw [arrow] (t3.east) -- ++(0.3,0) |- (eval.west);
				\draw [arrow] (t4.east) -- ++(0.3,0) |- (eval.west);
								
				% Нижняя петля обратной связи (проходит оптимально под блоком t4)
				\draw [arrow] (eval.south) -- ++(0,-2) -| node[pos=0.25, below, font=\scriptsize] {Не сошлось: модификация вектора параметров $\theta$ и запуск новой итерации} (algo.south);
				
			\end{tikzpicture}
		}
	\end{center}
\end{frame}